% GUID v6, Part III: one count of every scored confrontation. Numbers: docs/simulation/tally.py, which reads the result files of
% predictions_vs_data.py, standard_model.py and exotic_hadrons.py and writes docs/simulation/results/tally.json (6 of 6 checks).
% Rerun tally.py after any of the three scripts; it also picks up the rows of Hypothesis H9 once exotic_hadrons.json carries them
% (top-level key "H9"), and the table below must then be updated from tally.json.
\section{The tally}\label{sec:tally}

This section counts the scored confrontations of Sections~\ref{sec:data}, \ref{sec:sm} and~\ref{sec:exotics} together, in the classes of
Section~\ref{sec:scoring}.

\begin{table}[h]
\centering\small
\caption{Every scored confrontation of Part~III. Shared: the prediction is that of general relativity, $\Lambda$CDM, the Standard Model
(QCD included) or quantum mechanics. GUID-specific rows are sensitive, or not yet sensitive (Section~\ref{sec:scoring}). In brackets: FAILs.}
\label{tab:tally}
\begin{tabular}{@{}lrrrrrr@{}}
\toprule
section & scored & OK & FAIL & shared & GUID-specific, & GUID-specific, \\
 & & & & & sensitive & not yet sensitive \\
\midrule
data (Section~\ref{sec:data}) & 31 & 28 & 3 & 18 (3) & 0 & 13 (0) \\
Standard Model (Section~\ref{sec:sm}) & 30 & 26 & 4 & 21 (3) & 1 (1) & 8 (0) \\
hadrons (Section~\ref{sec:exotics}) & 9 & 4 & 5 & 9 (5) & 0 & 0 \\
\midrule
all & 70 & 58 & 12 & 48 (11) & 1 (1) & 21 (0) \\
\bottomrule
\end{tabular}
\end{table}

Table~\ref{tab:tally} gives the count \chk{tally:total}. Of 70 scored confrontations, 58 hold at 95\% and 12 fail. The conditional rows of the
plateau potential (one FAIL), the consistency rows and the INFO rows are not counted. Rotation curves and the remnant
(Sections~\ref{sec:sparc} and~\ref{sec:remnant}) compare models by $\Delta$BIC and score no rows, so they are not counted either.

\paragraph{The FAILs.}
\begin{itemize}[nosep,leftmargin=*]
\item \emph{Shared, eleven.} Three are tensions of $\Lambda$CDM that share data: evolving dark energy, the Hubble constant and curvature
(D3.3, D3.4, D3.7). Three are tensions of the Standard Model: the CDF W mass (SM2d.4) and the electron's $g-2$ with $\alpha$ from rubidium
and from caesium (SM3.8, SM3.9), which are one disagreement between the two values of $\alpha$. Five are masses of a colour-spin model on
compact trees: $X_0(2900)$, $T_{c\bar s0}(2900)$, $X(6900)$ read as $2^{++}$, $P_c(4312)$ and $P_{\psi s}(4338)$.
\item \emph{GUID-specific and sensitive, one.} The preregistered size of $R=\dd\ln\mu/\dd\ln\alpha$ (SM3.1) is refuted by derivation, not by a
datum.
\item \emph{GUID-specific and not yet sensitive, none.}
\end{itemize}

\paragraph{The classes.} The thirteen GUID-specific data rows test the information scalar (D1) and Bridge~B (D6). The information scalar's
departures scale with a coupling that may be zero, Bridge~B's decoherence lies a factor $4\times10^{12}$ or more below what the largest
superpositions could see, and its diffusion lies deep inside the allowed window. Eight Standard Model rows are GUID-specific but not yet sensitive. Dirac neutrinos give no double-beta signal, but
Majorana neutrinos can give none either (SM2b.1, SM2b.2). Decoupled right-handed partners change $N_{\rm eff}$ negligibly (SM2b.4, SM2b.5).
The scalar's $\mu$--$\alpha$ line and its share of $g-2$ scale with the free coupling (SM3.2, SM3.4 to SM3.6). The hadron rows test a
standard colour-spin model on the compact trees of Hypothesis~\ref{h:tree}. That hypothesis cannot fail on a single state, and the compact picture
predates GUID, so they are counted as shared. Table~\ref{tab:tally} would change if a reader classed them as GUID-specific: the five
model FAILs would then be GUID-specific and sensitive.

\paragraph{What the count says.} Forty-eight rows test what GUID adopts from established physics and carry eleven of the twelve FAILs.
Twenty-two rows test what GUID adds. Of these, one is refuted by derivation, and twenty-one hold but are not yet sensitive to
GUID's departures. No scored datum yet reaches where GUID departs from the theories it adopts. The class of each row and its reason
are in \texttt{tally.json} (6 of 6 checks).
